\documentclass[10pt, a4paper]{memoir}

\usepackage[utf8]{inputenc}
\usepackage[T1]{fontenc}
\usepackage[polish]{babel}

\usepackage{graphicx}
\usepackage{caption}
\usepackage{verse}
\usepackage{verse}
\setlength{\vleftmargin}{0pt} % Zeruje lewy margines pakietu verse
\indentpattern{0}             % Wymusza brak wcięć dla kolejnych linijek

\pagestyle{plain}
\linespread{1.3}

\setsecnumdepth{none}
\setlength{\leftmargini}{0pt}

\title{Listy do S.}
\author{Kate Malec}
\date{\the\year}

\begin{document}

\maketitle
\thispagestyle{empty}
\newpage

\thispagestyle{empty}
\vspace*{\fill}
\begin{flushright}
    \itshape
    Z gorącymi pozdrowieniami
\end{flushright}
\vfill

\newpage
\tableofcontents
\newpage

% --- WIERSZE ---

\part{Listy obecne}

\section{Koronacja}
\begin{verse}
Mój drogi\\
próbowałam na wiele sposobów,\\
zagryźć,\\
zapić,\\
wyprzeć,\\
zapomnieć,\\
lecz nie mogę.\\
Zbyt głęboko udało Ci się wejść mi~w~głowę.

W ciszy samotnie\\
Raz po raz\\
palcem sztywnym dotykam\\
za skronią, gdzie rozsądek\\
kiedyś jeszcze zipał,\\
teraz jest pusto.

By z braku nie oszaleć\\
towarzystwo Twoje\\
zepchnęłam na margines\\
rzeczywistości,\\
etykietkę Ci na czoło nakleiłam\\
\textit{,,wytwory wyobrażeń''},\\
nowym jesteś Hannibalem.\\
Na zawsze uwięziony\\
w splotach siedzisz mojej głowy.

Po części stało się dokładnie,\\
jak mówiłeś:\\
będziesz mnie uwierał,\\
że osobno --- przeoczyłeś.
\end{verse}

\newpage
\section{Czas wspólny}
\begin{verse}
Nie jestem w stanie wyrazić prostymi słowami,\\
co ze mną zrobił czas wspólny między nami.

Uszczerbki na zdrowiu miewam rozmaite:\\
od bólów, po zdrowie psychiczne rozmyte.

Żołądek, jelita odmówiły pracy.\\
Staram się leczyć, może ciało mi wybaczy.

Mimo wszystko nie powiem, że żałuję,\\
choć noc w noc serce tłukące się tulę.

Teraz tylko czekam aż dojdę do siebie\\
i będę w stanie w podzięce coś zrobić dla Ciebie.

Nie jestem naiwna, nie powiem \textit{miłość} skrycie,\\
lecz na nic próbować zagryźć to przeżycie.

Myślę, że Ty podobne miewasz dygresje,\\
Na papierze i płótnie spotkamy się jeszcze.
\end{verse}

\part{Stare wiersze}
\thispagestyle{empty}

\newpage
\section{Porcelanowe ptaki}
\begin{verse}
Ileż to krów, Hermesie, znów skradłeś Apollowi?\\
Tu słońce kolor niesie, śmierć berkiem kwiaty goni.\\
Próżno radość zakwitła, chwila palce ominie,\\
Nie złapiesz sztywną ręką, cień je rychło opłynie.

Myśli strugają chaos, tony dzielą na czworo,\\
To porcelanowe ptaki, puszczają do mnie oko.\\
Słucham leniwych dźwięków, świata śniegiem głuchego,\\
Odliczam dni strumieniem z kropli humoru wątłego.

Liczę nasze uśmiechy, wzajemne, choć osobno.\\
Nie obawiaj się, drogi, utrzymam dłoń łagodną.
\end{verse}

\newpage
\section{Hermès}
\begin{verse}
Pewien uroczy dżentelmen razu pewnego\\
Wziął mnie w rejs gondolą wieczoru zlewnego.\\
Świat razem zwiedzamy, urocze zakątki,\\
Umysłu bądź ciała --- plączą mi się wątki.

Oniemiała z wrażenia, zdjęta romantyzmem\\
Nie pomnę czym rozsądek, gnana optymizmem.\\
Wtem, w pół drogi będąc, spoglądam przez ramię,\\
Drogi! Toż my wody Styksu przemierzamy śmiale!

Droga, rzecze --- Spokojnie, Hermès moje imię,\\
Na brzeg rzeki cię później odprawię szczęśliwie.\\
Płyniemy więc dalej, niepomna rozsądku,\\
Ja, biedna konkubina, w poplątanym wątku.

Toniemy w radości, lecz czas bezlitosny\\
Pogania nam dłonie w uścisku miłosnym.\\
Tupie nogą zegar, czas wciągnąć ubrania,\\
Wybił koniec wspólnego Styksu zwiedzania.

Dobija do brzegu, bezpiecznie sprowadza,\\
Macham mu, żegna się, grunt żartu nie zdradza.\\
Odwracam się nagle i z lękiem spostrzegam:\\
To zły brzeg rzeki! W ufności poległam.

Odwracam się, wołam, każę mu zawracać,\\
Lecz próżno go szukać, ślad po nim się zatarł.\\
W kieszeni mały liścik --- ,,Znudził mi się zapach''.\\
Romantyzm już chyba przestał się opłacać.

Teraz stojąc boso na złym Styksu brzegu\\
Przeklinam tak ufność, jak zapach lubego,\\
Liczę włosy siwe, tulę duszę martwą,\\
Myślę: jak pachnie teraz? i czy było warto?
\end{verse}

\newpage
\section{Oda do Zimnej Marleny}
\begin{verse}
Marleno, Marleno,\\
posuń się, kochana,\\
jak za dawnych, ciepłych nocy\\
poleżymy znów do rana.

Pod kołdrą z mokrej ziemi,\\
pod foliowym worem,\\
na gorsze czasy ciebie\\
zakopałem tu, w ogrodzie.

Marleno, moja sztywna,\\
przesuńże kolana,\\
za życia uwierałaś,\\
po śmierci żadna zmiana.

Marleno, moja ciepła,\\
Gdy od środka gnijesz,\\
W fiolecie ci do twarzy,\\
Miłości aura bije.

Grzechoczą twoje kości,\\
stawy skrzypią twardo\\
pomyślałby ktoś jeszcze,\\
że śmiejesz się raźno.

Marleno, moja zimna,\\
za życia sztywną byłaś,\\
tyle w tym dobrego:\\
po śmierci niewiele się zmieniłaś.
\end{verse}

\newpage
\section{Konkubinat}
\begin{verse}
--- A Pan, drogi Prezesie, gdzie spędza te wakacje?\\
--- Ja?\dots --- Prezes pyta skromnie --- \dots a, w konkubinacie\dots
\end{verse}

\newpage
\section{Jałowy step}
\begin{verse}
Bosą nogą idź przez step,\\
Gdzie horyzont stracił końce.\\
Samotnością wypij wiatr,\\
duszę, która nie zna ojca.

Krok ciężki dusi morze traw,\\
ciąg martwych pustych dat.\\
Zostaw lęk, że ktoś nie przyjdzie,\\
w lęku złoto straci smak.

Więź swobody skórę zdziera\\
lękiem splata troski warkocz.\\
Wolność pozostaje pustką,\\
nie poczeka nikt na warcie.

Lepszy głód nie bez imienia\\
strach ubrany w życia zbroję.\\
Lepszy płacz, gdy kwitnie zboże,\\
Niż pusta przestrzeń samotności.
\end{verse}

\newpage
\section{Pokuta}
\begin{verse}
Generał, co utracił głos\\
w traw morzu nago stoi\\
Nałykał się już w życiu dość\\
Pieprzu oraz soli.

Życie --- struna, mocno drga\\
Żałośnie skrzypią drzwi\\
Sandały trzeszczą, płynie w dal\\
ulotnej chwili bzyk.

Z nosa sączy krew kroplami,\\
pomyślisz, że to łzy.\\
To tylko w ciemię bity głąb\\
pokutą karmi żwir.
\end{verse}

\end{document}